% =====================================================================
% Secao de Implementation (nova). Descreve o sistema concreto, fundamentada
% no codigo do repositorio. Idioma: ingles. Sugestao: inserir logo apos a
% Secao "DNAT Architecture" e antes da "Security Analysis".
% =====================================================================

\section{Implementation}
\label{sec:implementation}

This section describes how the three planes are realized in the prototype, so
that the security properties evaluated later can be traced to concrete mechanisms
rather than to design intent alone. The system is deployed as three containers on
a shared runtime---\texttt{dnat-client} (CVM1), \texttt{dnat-builder} (CVM2), and
\texttt{dnat-executor} (CVM3)---and is contrasted with a single-container
baseline that colocates the same logic.

\subsection{Control Plane (CVM1)}

The control plane hosts the marketplace API, a local IPFS node, and the
blockchain client. Application and dataset registration, access checks, and
execution orchestration are exposed through a small HTTP API. On application
registration, CVM1 does not install dependencies itself: it assembles a build
bundle (the source file plus a manifest of declared dependencies) and forwards it
to the builder. When the builder returns an \texttt{application.ext4} artifact,
CVM1 encrypts it before publication. The encryption uses AES-256-GCM over a
gzip-compressed artifact, with a versioned binary envelope (magic header, length
prefix, JSON header carrying the IV, authentication tag, and metadata). Only the
resulting encrypted blob and a manifest are published to IPFS; the plaintext
artifact never leaves the control plane. Access control is enforced on chain: an
execution request first calls \texttt{hasAccessByIds}, and the workflow aborts
with \emph{access denied} if the caller has not purchased access to both the
dataset and the application.

\subsection{Build Plane (CVM2)}

The builder exposes a single \texttt{/build} endpoint that accepts a bundle and
launches an ephemeral Firecracker microVM to resolve dependencies. Unlike the
executor, the build microVM is given a network interface, because resolving and
compiling Python wheels may require egress to package indexes. Egress is
constrained on the host side: a TAP device is created and iptables rules
\texttt{REJECT} traffic from the guest toward private ranges
(\texttt{10.0.0.0/8}, \texttt{172.16.0.0/12}, \texttt{192.168.0.0/16},
\texttt{127.0.0.0/8}, \texttt{169.254.0.0/16}, and the host TAP address) while
allowing the remaining outbound traffic via masquerading. Inside the guest, an
artifact builder produces a read-only \texttt{ext4} image containing the
application source under \texttt{/app}, its \texttt{site-packages}, a lock file of
resolved dependencies, and a marker file used later for disk discovery.
Dependencies are first wheel-built and then installed with \texttt{--no-index}
from the local wheelhouse, and reusable wheels are persisted in a size-pruned
cache that is the builder's only durable state. After the build, the host returns
only the artifact and build metadata; the per-request build state is discarded.

\subsection{Execution Plane (CVM3)}

The executor exposes a single \texttt{/execute} endpoint. For each request it
copies a fresh root overlay, creates new input and output \texttt{ext4} disks,
optionally attaches the application artifact disk read-only, and launches a
Firecracker microVM with one vCPU and 512~MiB of memory and---crucially---no
network interface. The guest does not receive disk roles out of band; it
discovers them by scanning candidate block devices for marker files, then
extracts the workspace and runs \texttt{workspace/run.sh}. The launcher script
tries a small, fixed set of invocation strategies so that applications can accept
the dataset either as \texttt{--dataset data/dataset.csv} or as a positional
argument; standard output and error are captured and persisted to the output
disk. Synchronization between host and guest uses the serial console: the guest
writes an \texttt{EXECUTION\_COMPLETE} marker, after which the host issues a
graceful shutdown, escalating to termination if needed. The host then mounts the
output disk and derives the HTTP response solely from the persisted result. A
\texttt{cleanup} trap unmounts every loop device and removes the entire per-run
working directory on exit, which is what makes execution state ephemeral.

\subsection{Asset Integrity and the Baseline}

Assets are bound to the chain by content hash: registration records the SHA-256
of the artifact, which lets a consumer detect tampering between IPFS and the
contract. The baseline used for comparison
(\texttt{docker/baseline-vm.compose.yaml}) runs the identical control, build, and
execution logic, including Firecracker isolation for the build and execution
microVMs, but inside a single container whose filesystem holds the IPFS
repository, the wheel cache, and the executions store, and whose environment
holds the asset-encryption key and the wallet private key. The only structural
difference from the compartmentalized deployment is therefore the colocation of
these assets, which is exactly the variable the security evaluation isolates.
